% ============================================================
% sec/03methodology.tex
% Required: problem formalization, pipeline, experimental protocol
% (splits, validation, metrics + justification), baseline plan,
% risks with mitigation.
% ============================================================
\section{Proposed Methodology}
\label{sec:methodology}

\subsection{Problem Formalization}
\label{subsec:formalization}
\todo{Mathematical formulation. Define the input space, the output
space, the objective function, and the constraints. For Track C, point
to the MDP in Section~\ref{sec:trackanalysis} and state the optimization
objective here.}

\begin{equation}
\label{eq:objective}
\todo{objective function}
\end{equation}

\subsection{Proposed Pipeline}
\label{subsec:pipeline}
\todo{End-to-end description: data ingestion, preprocessing, model,
training loop, evaluation, and the interface the Stage 3 agent will
consume.}

\begin{figure}[!t]
  \centering
  \safeimage[width=\columnwidth]{pipeline.png}
  \caption{Proposed pipeline. \todo{caption}}
  \label{fig:pipeline}
\end{figure}

\subsection{Experimental Protocol}
\label{subsec:protocol}

\subsubsection{Data Partitioning}
\todo{Train / validation / test split with exact proportions. State
whether the split is random, stratified, temporal or by group, and
justify the choice against the alternatives.}

\subsubsection{Validation Strategy}
\todo{Holdout, k-fold, nested CV, or episodic evaluation over N seeds.
Justify. State how hyperparameters are selected without touching test.}

\subsubsection{Metrics}
\begin{table}[!t]
\centering
\caption{Evaluation metrics and their justification.}
\label{tab:metrics}
\footnotesize
\begin{tabularx}{\columnwidth}{@{}l Y@{}}
\toprule
\textbf{Metric} & \textbf{Why this metric, and what it beats}\\
\midrule
\todo{primary}   & \todo{justification}\\
\todo{secondary} & \todo{justification}\\
\todo{secondary} & \todo{justification}\\
\bottomrule
\end{tabularx}
\end{table}

\subsubsection{Data Leakage Prevention}
\todo{Name every leakage vector considered (target leakage, test set
touched during tuning, temporal leakage, group leakage across splits)
and the concrete measure taken against each one.}

\subsection{Baseline Plan}
\label{subsec:baselines}
\begin{enumerate}
  \item \textbf{Trivial baseline.} \todo{e.g. majority class, random
        policy, rule-based controller.}
  \item \textbf{Competitive classical baseline.} \todo{The strongest
        non-deep or non-RL method reasonably applicable.}
  \item \textbf{Literature baseline.} The published target declared in
        Section~\ref{subsec:litbaseline}.
\end{enumerate}

\subsection{Reproducibility Contract}
\label{subsec:reproducibility}
\todo{Fixed seeds and where they live, saved configuration files, run
logging, hardware used, dependency pinning, and the exact command that
reproduces each reported number.}

\subsection{Risks and Mitigation}
\label{subsec:risks}
\begin{table}[!t]
\centering
\caption{Identified risks and their mitigation.}
\label{tab:risks}
\footnotesize
\begin{tabularx}{\columnwidth}{@{}P{2.0cm} c Y@{}}
\toprule
\textbf{Risk} & \textbf{Impact} & \textbf{Mitigation / trigger}\\
\midrule
\todo{risk} & \todo{H/M/L} & \todo{mitigation}\\
\todo{risk} & \todo{H/M/L} & \todo{mitigation}\\
\todo{risk} & \todo{H/M/L} & \todo{mitigation}\\
\bottomrule
\end{tabularx}
\end{table}

\subsection{Planned Agent Integration}
\label{subsec:agent}
\todo{How the trained model will be consumed by an intelligent agent in
Stage 3: LLM with tool calling, RAG system, multi-agent architecture, or
an RL agent operating in its environment. State the interface and how the
consumption will be verified.}